\section{Mobil és web standardok (reszponzív dizájn)}
Az okostelefonok korszaka újratervezésre kényszerítette a webalkalmazásokat \cite{w3c2008mobileBestPractices}. Az olyan webes szabványok, mint a CSS3 media query-k és a flexibilis elrendezések, lehetővé tették a reszponzív dizájnt: olyan oldalak létrehozását, amelyek alkalmazkodnak a különböző képernyőméretekhez \cite{marcotte2010responsiveWebDesign,w3c2012mediaQueries}. A W3C HTML5 és CSS3 specifikációi támogatást nyújtottak mobilbarát funkciókhoz is (pl. érintőképernyő általi események) \cite{w3c2014html5,w3c2013touchEvents}.

Sok weboldal átvette a „mobile-first” tervezési szemléletet \cite{wroblewski2011mobileFirst}. A körülbelül 2015-ben megjelent Progressive Web Appok (PWA-k) Service Workereket és manifest fájlokat használnak annak érdekében, hogy natív alkalmazásokhoz hasonló élményt nyújtsanak (offline mód, telepíthető ikonok) \cite{russell2015progressiveApps,mdnWhatIsPwa}. Ezek a modern HTML5/JavaScript technológiákra és böngésző API-kra (kamera, helymeghatározás stb.) épülnek, hogy fejlett funkcionalitást biztosítsanak \cite{mdnGetUserMedia,w3c2024geolocation}.
